% Ciclo de desenvolvimento do RECLin-PT, com o estado de cada etapa conferido
% nos artefatos (branch resultados-colab, commit bb9457f) e no protocolo.
\begin{figure}[htbp]
  \centering
  \caption{Ciclo de desenvolvimento do RECLin-PT e estado atual de cada etapa.}
  \label{fig:ciclo}
  \begin{tikzpicture}[
      font=\sffamily\footnotesize,
      etapa/.style={rounded corners=3pt, minimum height=1.25cm, text width=3.0cm, align=center, inner sep=3pt},
      concluida/.style={etapa, draw=black!80, fill=black!20, line width=0.8pt},
      parcial/.style={etapa, draw=black!80, line width=0.8pt, fill=white, pattern=north east lines, pattern color=black!25},
      pendente/.style={etapa, draw=black!60, dashed, line width=0.8pt, fill=white},
      seta/.style={-{Stealth[length=2.2mm]}, line width=0.8pt},
      estado/.style={font=\sffamily\scriptsize\itshape, text=black!70, below=1.5pt}
    ]
    % linha 1 (da esquerda para a direita)
    \node[concluida] (e1) at (0,0)    {\textbf{1.} Preparação\\ dos dados};
    \node[concluida] (e2) at (3.75,0) {\textbf{2.} Geração de\\ pares candidatos};
    \node[parcial]   (e3) at (7.5,0)  {\textbf{3.} \emph{Fine-tuning}\\ (\emph{encoder} $\times$ estratégia)};
    \node[concluida] (e4) at (11.25,0){\textbf{4.} Seleção de épocas\\ e limiares no DEV};
    % linha 2 (da direita para a esquerda)
    \node[parcial]   (e5) at (11.25,-3.4) {\textbf{5.} Comparação de\\ estratégias e \emph{encoders}};
    \node[pendente]  (e6) at (7.5,-3.4)   {\textbf{6.} Congelamento\\ da configuração};
    \node[pendente]  (e7) at (3.75,-3.4)  {\textbf{7.} Avaliação única\\ no TEST};
    \node[pendente]  (e8) at (0,-3.4)     {\textbf{8.} Seleção do\\ modelo final};
    \draw[seta] (e1) -- (e2);
    \draw[seta] (e2) -- (e3);
    \draw[seta] (e3) -- (e4);
    \draw[seta] (e4.south) -- (e5.north);
    \draw[seta] (e5) -- (e6);
    \draw[seta] (e6) -- (e7);
    \draw[seta] (e7) -- (e8);
    % retorno 5 -> 3
    \draw[seta, dashed] (9.95,-2.775) -- (9.95,-1.7) -- (8.35,-1.7) -- (8.35,-0.625);
    \node[font=\sffamily\scriptsize, anchor=north east] at (9.85,-1.78) {novos \emph{encoders} ou estratégias};
    % estados
    \node[estado] at (e1.south) {concluída};
    \node[estado] at (e2.south) {concluída};
    \node[estado, xshift=-6mm] at (e3.south) {parcial$^{a}$};
    \node[estado, xshift=-8mm] at (e4.south) {concluída$^{b}$};
    \node[estado, xshift=-8mm] at (e5.south) {parcial$^{c}$};
    \node[estado] at (e6.south) {pendente};
    \node[estado] at (e7.south) {pendente};
    \node[estado] at (e8.south) {pendente};
    % legenda
    \begin{scope}[shift={(3.9,-4.95)}, font=\sffamily\scriptsize]
      \node[draw=black!80, fill=black!20, minimum width=5mm, minimum height=3.2mm, rounded corners=1pt] (l1) at (0,0) {};
      \node[right=1mm of l1] (t1) {concluída};
      \node[draw=black!80, fill=white, pattern=north east lines, pattern color=black!25, minimum width=5mm, minimum height=3.2mm, rounded corners=1pt, right=6mm of t1] (l2) {};
      \node[right=1mm of l2] (t2) {parcial};
      \node[draw=black!60, dashed, fill=white, minimum width=5mm, minimum height=3.2mm, rounded corners=1pt, right=6mm of t2] (l3) {};
      \node[right=1mm of l3] {pendente};
    \end{scope}
  \end{tikzpicture}\\[4pt]
  {\footnotesize $^{a}$ O BioBERTpt-all foi ajustado em todas as estratégias treinadas; o BERTimbau-base, apenas no baseline; o BioBERTpt-clin e o \emph{fine-tuning} restrito não têm execução. $^{b}$ Para as execuções existentes e verificadas. $^{c}$ Estratégias comparadas com o BioBERTpt-all; \emph{encoders}, apenas na estratégia baseline; comparações descritivas no DEV.}\\[2pt]
  {\footnotesize Fonte: elaborado pelo autor.}
\end{figure}
